\uuid{Jicq}
\titre{Application linéaire}
\auteur{hanani ; mzari}
\organisation{exo7}
\licence{CC BY-NC-SA 4.0}
\datecreate{2020-01-20}
\niveau{L1}
\module{Algèbre}
\chapitre{Application linéaire}
\sousChapitre{Définition}
\difficulte{1}
\theme{Applications linéaires, Application linéaire}
\exo7qcmid{lille-2.applications-lineaires.002}

\contenu{
\begin{question}
\qclassification{107.01}
\qtags{motcle=Application linéaire}
On considère les deux applications suivantes :
$$\begin{array}{rccc}f:&\Rr^2&\to&\Rr^2\\
& (x,y)&\to & (x,y^2) \end{array} \quad \mbox{et} \quad \begin{array}{rccc}g:&\Rr^2&\to&\Rr^2\\
& (x,y)&\to &(x,-x). \end{array}$$
Quelles sont les assertions vraies ?
\begin{answers}
\good{$f(0,2)=(0,4)$}.
\bad{$f$ est une application linéaire}.
\good{$g(0,0)=(0,0)$}.
\good{$g$ est une application linéaire}.
\end{answers}
\begin{explanations} L'application $f$ n'est pas linéaire. Contre-exemple : $f(2(0,1))=f(0,2)=(0,4)$ et $2f(0,1)=(0,2)$. On vérifie que l'application $g$ est linéaire.
\end{explanations}
\end{question}
}
